\documentclass[12pt]{article}
\usepackage[utf8]{inputenc}
\usepackage[T1]{fontenc}
\usepackage{hyperref}
\title{Wajuk Language Resources for K-12 Curriculum Development}
\author{Compiled by Claude Sonnet 5.5 (claude-sonnet-5-5) for Andr\'as Kornai}
\date{}
\begin{document}
\maketitle

\section{Introduction}
Wajuk (ISO 639-3 xwj; also spelled Whadjuk) is the Noongar dialect of the Perth area in Western Australia. ISO lists it as extinct, but it survives as part of the Noongar language, which has living speakers and an active revival. Speaker figures for Noongar are given as under 240 (2014) and 475 (2016 census). Most of the resources below cover Noongar generally, with some Perth-area focus. They are not Wajuk-only.

\section{Primary Educational Resources}
\begin{itemize}
\item Sharon Gregory, \emph{Noongar Language Dictionary}, a beginner's resource for students and adults in the Fremantle and Perth area (Koort-kadak Consultancy, first published 2014, updated 2019), as listed in a library catalogue: \url{https://cockburnlibrary.spydus.com/cgi-bin/spydus.exe/FULL/OPAC/BIBENQ/114780686/15510566,545}
\item A school library guide to Reconciliation Action Plan resources, including language: \url{https://libguides.allsaints.wa.edu.au/RAPresources/language}
\end{itemize}

\section{Linguistic and Reference Materials}
\begin{itemize}
\item Rose Whitehurst, \emph{Noongar Dictionary: Noongar to English and English to Noongar} (Noongar Language and Culture Centre, 1992), mentioned in the search summary; no direct link was seen.
\item Wikipedia, Noongar language (including Wajuk): \url{https://en.wikipedia.org/wiki/Noongar_language}
\item AIATSIS record W41, Noongar: \url{https://collection.aiatsis.gov.au/node/1510236}
\item Bookshop listing for \emph{Nyoongar Legacy: the naming of the land and the language of its people}: \url{https://elizabethsbookshop.com.au/shop/australiana/west-australiana/nyoongar-legacy-the-the-naming-of-the-land-and-the-language-of-its-people}
\end{itemize}

\section{Religious and Community Resources}
The search summary mentions the Noongar Boodjar Language Cultural Aboriginal Corporation as working on a dictionary and language database. This comes from a secondary travel site (\url{https://citiesinsider.com/country/australia/perth/indigenous-languages/en}) and should be verified with the organisation.

\section{Digital and Audio Resources}
No Wajuk-specific audio, video or software was found in these searches.

\section{Curriculum Development Context}
Wajuk is part of the living Noongar revival, so a heritage curriculum would be developed with Whadjuk and wider Noongar custodians. It could draw on existing Noongar dictionaries and school materials, adapting them to Whadjuk forms. Wilfrid Douglas's mid-twentieth-century work (grammar and dictionary) is a documented source; the Noongar Boodjar Language Centre is working on dictionaries and learning materials. Community permission is necessary for any use.
\end{document}